\section{Source-anchored discovery with the maintained search priorities}
\label{sec:anchored-producer}

The representation from Section~\ref{sec:anchored-replay} now also supports
native producer discovery. Previously, even consecutive raw primitive moves
rebuilt ordinary words after each round. The new producer keeps the original
roots and a monomial image table through candidate discovery, primitive-width
queries and subsequent contractions, then exports once. Its arithmetic is
separate from the independent checker. No checker file, certificate schema
or source prerequisite changes.

\subsection{Recomputing eligibility on the fixed source}

The producer evaluates signed occurrence counts on the surviving source
roots under the current image table. A terminal mapped to $z^k$ contributes
$|k|$ positive or negative occurrences according to both the original sign
and the sign of $k$. Concatenation adds counts. Each node retains at most two
current labels; a larger support is represented by an ineligible marker.
A current root is a candidate precisely when it has two labels and neither
label occurs with both signs. Slot order remains unchanged.

Crucially, these summaries are recomputed after each image update. A source
node with three original labels can become a valid two-label word after two
labels acquire the same image. Permanently caching its original ineligibility
would lose such donors. The old producer's immutable-node cache was valid
because substitution created new ordinary roots; source anchoring changes
that cache key.

The maintained projection and forest planners consume this current candidate
snapshot. Their only new internal parameter is a primitive-power callback.
It evaluates the source word's exact prefix-height extrema under the current
signed images, with no intermediate ordinary powering circuits. A coherent
word with an occurrence count of one retains the existing immediate primitive
case. Other candidates use the same gcd-normalized exponent vector and width
criterion as before. Forest component merging, child selection, disjoint-pair
selection, donor-slot order and the shared within-round power cache are
unchanged. Hence the producer discovers the same greedy move, not merely
some algebraically equivalent contraction.

The new producer composes disjoint images directly. For a forest it resolves
parent chains with memoized cumulative powers. This differs from the checker's
independent forest validation and breadth-first reconstruction. Producer
planning proves internal eligibility; only full source-bound replay can turn
the resulting trace into an unknot verdict.

\subsection{Preserving higher-priority operations}

The ordinary producer still discovers the first move of an epoch. If that
move leaves at most two generators, its existing direct updater is retained.
The same direct route is retained when every primitive-vector coordinate in
the move has absolute value one. Such a move only relabels signed terminals:
it creates no nontrivial binary powering circuit. Complete source timings
showed that importing these ordinary cases into lazy discovery added overhead.

Otherwise the lazy block starts from the pre-move ordinary roots and consumes
the already discovered move. Thus entry requires a nonunit monomial image
and room for further raw discovery. This is an adaptive implementation choice
between two exact producers; it does not suspend any donor search, change
its order, or introduce a new mathematical hypothesis. A block already entered
continues through unit moves as well. The ordinary unit path remains covered
by the earlier polynomial raw-epoch theorem. It repeats discovery until no primitive move
is available, rank reaches two, or a higher-priority singleton batch becomes
available. Returning at rank two preserves the caller's existing primitive
terminal test, including the option that disables that test.

When singleton batches are enabled, their priority is evaluated exactly on
the lazy state. Every source node carries raw length, a presence bitmask and
a repeated-label bitmask. At a terminal with image $z^k$, the repeated bit is
set exactly when $|k|>1$. At concatenation, repeated labels are the union of
both children's repeated sets and their support intersection. A reverse pass
propagates root multiplicities through source edges, including duplicate
children, and accumulates each label's exact global occurrence count weighted
by $|k|$. These summaries recover the maintained score
$(\ell-2)\,\mathrm{count}(g)$, length, slot and generator tie-breaks.
The existing acyclic singleton selection routine decides whether at least
two donors take precedence. If so, the block exports and returns to that
ordinary search branch; it does not silently execute a projection first.

A stall also exports before power-pair deletion, normalization, overlap moves
or Whitehead search. The caller may repeat a boundary eligibility query on
the exported state; this conservative handoff preserves the existing control
flow. Ordinary source reconstruction and final endpoint tests remain outside
the private image-table interface. Roots, live labels and emitted moves are
published only after all discovery and export work succeeds. Limits and
cancellation cannot publish a partial trace. As with the checker, speculative
ordinary export nodes may remain allocated on failure and stay charged.

\subsection{Encoded cost and the remaining global problem}

Use $r,m,S_0,B,\Lambda$ from Section~\ref{sec:anchored-replay}. The fixed-source
saturated-kernel theorem again bounds cumulative exponents by
$\Lambda=O(r(B+\log(r+1)))$. Each successful primitive round removes a
label, so a block has at most $r-1$ successful rounds and one final unsuccessful
query. Signed counts, lengths and height extrema have $O(B+\Lambda)$ bits.
No ordinary nodes are allocated by discovery or image composition; final
export gives
\[
 N_{\rm final}\le N_0+O(S_0+r\Lambda).
\]

Unlike supplied-certificate verification, discovery can test unsuccessful
donors. The number of successful contractions alone does not bound these
queries. A simple safe bound allows $m$ candidate roots in each of $O(r)$
rounds, each visiting at most $S_0$ source nodes. The implementation obtains
topological order by sorting reachable ordinary node handles, so one may
bound this portion by
\[
 O\!\left(r(m+1)(S_0+1)\log(S_0+2)+rm+r^2+r\Lambda\right)
\]
node, comparison and arithmetic operations, with integer bit costs separate.
This is deliberately a whole-block polynomial bound rather than a claim
that every failed profile is charged to a removed generator.

With singleton priority enabled, its source summary passes remain polynomial
in the same parameters. Each round supplies at most $mr$ singleton candidates
to the existing finite acyclic selector, with bitmasks of at most $r$ bits;
its sorting and reachability work must also be counted. There are at most
$r$ such calls in a block. These observations establish polynomial encoded
maintenance and discovery within this restricted source model. They neither
bound the number of nonmonomial source changes nor establish that these moves
find an unknot proof for every diagram. The general quasipolynomial objective
therefore remains open.

\subsection{Audit, adaptive correction and complete-call evidence}

The native \code{anchored\_producer.py} research driver compares
against \code{339c5359d34b}, which already contains the source-anchored checker.
It loads the predecessor in an isolated namespace and pins 449 current source
and fixture files before and after each run. Three pre-existing producer
modules change: the search dispatcher and the two greedy planners' internal
profile callbacks. The source-anchored producer is a new module. All independent
checker files, including the compressed dispatcher, remain byte-for-byte equal
to the baseline.

The final maintained suite passes all 1,137 tests in 219.992 seconds. Focused
tests compare exact greedy projection, forest and singleton-priority choices
on 240 randomized signed presentations, including unrelated and empty roots.
Other cases cover binary powers, checker-disabled production, producer-disabled
source replay, unit-versus-growth dispatch, node/work exhaustion and cancellation
without partial public state or trace publication.

The source audit covers 102 diagrams in six modes: default, projection,
forest, combined singleton/projection, combined singleton/forest, and projection
with the rank-two terminal disabled. Both packages return identical complete
certificates or the same absence of a certificate. The two combined Gordian
modes exhaust the same five-million-work allowance in both engines; these
are recorded incomplete discovery attempts, not successes. Every positive
certificate is replayed by both engines in both literal and compressed modes.
The adaptive rule leaves these ordinary diagram cases on the direct producer.

Six further cases start from actual five- and seventeen-strand stabilized
circles and a supplied legal Whitehead power with exponent $2^b$, for
$b=4,128,512$. Subsequent primitive discovery is automatic and activates one
two-round lazy block in every case. Complete certificates, including the
exposure prefix, agree with the predecessor and pass both compressed source
checkers; the $b=4$ cases also pass both literal checkers. The full audit totals
1,328 source replays and takes 34.060 seconds while overlapping the suite;
its duration is not a performance measurement. On the five-strand $b=512$
case, producer work falls from 54,305 to 39,465 and nodes from 2,591 to 2,078.
On the seventeen-strand $b=4$ case, work instead rises from 5,099 to 5,243:
the adaptive criterion is a practical choice, not a proof that every accepted
entry is faster. Automatic discovery of the exposure prefix is not claimed.

Abstract infinite-cyclic capacity cases additionally compare every move,
terminal claim and retained run-length word. At rank 32 and 128 local bits,
the balanced case changes final nodes from 8,170 to 4,250 and total producer
work from 160,464 to 93,032. The 32-generator star has 385 nodes in both
engines, with work 109,465 versus 56,429. A 16-generator, 512-bit chain in
projection mode changes nodes from 38,598 to 29,316 and work from 568,655
to 447,260. The same chain in forest mode ends in one batch and retains
exactly the old route: 22,635 nodes and 208,807 work. These are algebraic
capacity checks, not new diagram-derived recognition successes.

\paragraph{Why entry became adaptive.}
The first serial timing run used lazy discovery whenever a move left more
than two labels. On a 129-strand circle, complete projection discovery and
replay changed from 143.994 to 194.017 ms, with paired old/new ratio $0.730$;
the forest case changed from 92.595 to 142.539 ms, ratio $0.664$. Both are
real regressions despite passing correctness tests. Signed unit images have
no nontrivial powers to amortize the repeated source-summary work. The final
entry rule therefore retains the cached ordinary path on those moves.
The complete initial audit, tests and all three benchmark runs remain in
\path{data/anchored-producer-before-gate-*}; they are not substituted for
measurements of the final implementation.

\paragraph{Final serial measurements.}
After the adaptive correction, the suite, audit and second code publication,
the driver reruns all three timing modes serially. Five shuffled rounds retain
both engines and an A/A copy of each. All 700 measured calls and 140 warm-ups
complete. The tables contain median times and median within-round ratios;
these are different estimators, especially when machine timings drift across
rounds. No incomplete baseline is used as the denominator of a speedup claim.

For algebraic stages, each timer includes fresh grammar construction, complete
greedy discovery, fresh source-grammar reconstruction and independent replay
through the rank-one endpoint. Exact output comparisons and proof serialization
are outside the timer. All words, moves and terminal claims agree. The five
cases take 21.687 seconds for 100 measured calls and 20 warm-ups.

\begin{center}\small
\begin{tabular}{lrrrrr}
\toprule
Case & Old ms & New ms & Old/new & Old A/A & New A/A \\
\midrule
star-4-0 & 0.460 & 0.508 & 0.885 & 1.021 & 1.117 \\
balanced-32-64 & 52.760 & 38.859 & 1.363 & 0.992 & 1.004 \\
star-32-128 & 63.602 & 52.502 & 1.233 & 0.990 & 0.974 \\
chain-16-512 & 474.307 & 413.783 & 1.143 & 0.983 & 1.009 \\
chain-16-512-forest & 286.595 & 284.324 & 0.993 & 0.992 & 0.993 \\
\bottomrule
\end{tabular}
\end{center}


The balanced, star and chain binary-growth cases give paired old/new ratios
$1.363$, $1.233$ and $1.143$. These are complete discovery-and-replay operations,
not the maintenance-only schedules timed in the preceding section. They also
use a different baseline, which already has the new checker; ratios from the
two sections must not be multiplied into an unmeasured combined speedup.
The forest chain ends in one batch and changes neither algorithm nor proof;
its ratio $0.993$ is a control. The sub-millisecond unit star has a new A/A
ratio $1.117$, so its $0.885$ cross-engine ratio is too noisy to establish a
reliable regression or gain.

The source-stage timers include fresh PD validation, discovery and mandatory
independent presentation recovery, full proof replay and endpoint verification.
Cases named ``exposed'' additionally construct the stated supplied Whitehead
prefix inside the timer. Choosing that prefix is not automatic. The eleven
cases take 14.103 seconds for 220 measured calls and 44 warm-ups.

\begin{center}\small
\begin{tabular}{lrrrrr}
\toprule
Case & Old ms & New ms & Old/new & Old A/A & New A/A \\
\midrule
circle-5-projection & 0.687 & 0.694 & 0.996 & 0.947 & 0.979 \\
circle-5-forest & 0.640 & 0.638 & 1.006 & 0.987 & 0.993 \\
circle-17-projection & 4.443 & 4.478 & 1.000 & 1.003 & 1.004 \\
circle-17-forest & 3.715 & 3.661 & 1.014 & 1.022 & 0.989 \\
circle-65-projection & 40.742 & 39.652 & 1.030 & 1.019 & 0.990 \\
circle-65-forest & 28.087 & 27.816 & 1.014 & 1.017 & 1.001 \\
circle-129-projection & 257.665 & 221.376 & 0.994 & 1.011 & 1.009 \\
circle-129-forest & 168.659 & 169.456 & 0.989 & 0.987 & 1.012 \\
circle-5-exposed-128 & 14.690 & 18.173 & 1.021 & 1.003 & 0.990 \\
circle-17-exposed-128 & 29.194 & 28.547 & 1.025 & 1.003 & 0.995 \\
circle-17-exposed-512 & 49.051 & 47.322 & 1.050 & 1.004 & 1.023 \\
\bottomrule
\end{tabular}
\end{center}


The 129-strand ordinary projection and forest ratios are now $0.994$ and
$0.989$, respectively. The large regression from unconditional entry is absent;
these cases retain ordinary production. The exposed cases give $1.021$,
$1.025$ and $1.050$. Their modest effects should be interpreted with the raw
samples. In the 129-strand projection run, old-engine times move from about
140 ms in the first round to about 258 ms in later rounds. In the five-strand
exposed case, the unpaired median time increases while the median paired ratio
is above one. Neither should be presented as a robust source-wide improvement.
The seventeen-strand, 512-bit exposed case has ratios above one in all five
paired rounds, but its new A/A median is $1.023$: the observed benefit remains
small compared with the algebraic capacity gains.

Finally, the whole-recognition control uses the nineteen-case corpus with
primitive projections and compressed search enabled in both engines. It
includes discovery and required source replay, preserving the same
20-million group-work, ten-second group, twelve-second total and 50,000-object
limits. All 380 measured calls and 76 warm-ups complete in 35.270 seconds.

\begin{center}\small
\begin{tabular}{lrrrrr}
\toprule
Case & Old ms & New ms & Old/new & Old A/A & New A/A \\
\midrule
survivor-00 & 15.166 & 15.293 & 0.996 & 1.000 & 1.001 \\
survivor-01 & 13.022 & 13.012 & 1.001 & 0.995 & 1.005 \\
survivor-02 & 14.241 & 14.243 & 0.983 & 1.005 & 1.004 \\
survivor-03 & 10.197 & 10.098 & 1.008 & 1.012 & 1.004 \\
mirror-03 & 22.814 & 23.176 & 0.992 & 0.996 & 0.980 \\
survivor-04 & 8.680 & 8.786 & 0.992 & 1.001 & 0.996 \\
survivor-05 & 17.662 & 17.725 & 0.991 & 0.993 & 1.000 \\
survivor-06 & 7.000 & 7.129 & 0.985 & 1.002 & 1.007 \\
survivor-07 & 18.588 & 18.693 & 1.003 & 1.003 & 0.991 \\
survivor-08 & 6.875 & 7.002 & 0.978 & 0.990 & 1.014 \\
mirror-08 & 19.171 & 19.303 & 0.996 & 0.993 & 0.967 \\
survivor-09 & 7.081 & 7.036 & 1.002 & 1.002 & 0.983 \\
survivor-10 & 6.216 & 6.064 & 1.000 & 1.012 & 1.008 \\
survivor-11 & 6.792 & 6.624 & 1.024 & 1.032 & 0.984 \\
gordian & 1311.242 & 1232.129 & 1.064 & 1.076 & 1.001 \\
trefoil & 0.056 & 0.055 & 1.011 & 1.015 & 1.018 \\
figure\_eight & 0.061 & 0.060 & 1.017 & 1.007 & 0.949 \\
conway & 0.989 & 0.993 & 0.996 & 1.000 & 1.002 \\
kinoshita\_terasaka & 1.030 & 1.028 & 0.985 & 0.998 & 0.998 \\
\bottomrule
\end{tabular}
\end{center}


Paired old/new ratios range from $0.978$ to $1.064$. Old A/A ranges from
$0.990$ to $1.076$ and new A/A from $0.949$ to $1.018$. In particular the
Gordian cross-engine ratio $1.064$ is smaller than its old A/A ratio $1.076$,
so it is not evidence for a recognition improvement. The source audit finds
no adaptive lazy entry on the ordinary corpus. Recognition coverage and
complete certificates are unchanged. These measurements support a targeted
representation choice and polynomial fixed-source discovery, not a general
recognition speedup or a quasipolynomial completeness theorem.

The final raw records and their table generator are:
\begin{flushleft}\small
\path{data/anchored-producer-kernels.json}\\
\path{data/anchored-producer-source.json}\\
\path{data/anchored-producer-pipeline.json}\\
\path{data/anchored_producer_tables.py}
\end{flushleft}
Each record retains every sample, source description, proof reference and
source pin. The read-only upstream
check finds \path{openai/math} unchanged at \code{fd4aeeb2ee4f}; no new external
mathematics is imported in this step.


